% Motor: XeLaTeX, fijado en .latexmkrc y en la receta del editor.
\documentclass[aspectratio=169,10pt,c]{beamer}

% !TEX root = slides.tex
% Preambulo comun: estilo UTP, colores, fuentes, plantillas de beamer y
% comandos de apoyo didactico. Se carga desde slides.tex tras \documentclass.

% La opción "c" de beamer reparte el espacio libre 1:1.5 (arriba:abajo), por lo
% que el contenido queda por encima del centro. Igualamos ambos rellenos.
\makeatletter
\define@key{beamerframe}{c}[true]{%
  \beamer@frametopskip=0pt plus 1fill\relax%
  \beamer@framebottomskip=0pt plus 1fill\relax%
  \beamer@frametopskipautobreak=0pt plus .5\paperheight\relax%
  \beamer@framebottomskipautobreak=0pt plus .5\paperheight\relax%
}
\makeatother

\usepackage{graphicx}
\usepackage{tikz}
\usepackage{xcolor}
\usepackage{iftex}

\ifPDFTeX
  \usepackage[utf8]{inputenc}
  \usepackage[T1]{fontenc}
  \usepackage{lmodern}          % tipografia escalable para el respaldo con pdflatex
  \renewcommand{\familydefault}{\sfdefault}
\else
  \usepackage{fontspec}
  \IfFontExistsTF{Arial}
    {\setsansfont{Arial}[BoldFont={Arial Bold}]}
    {\setsansfont{Helvetica}[BoldFont={Helvetica Bold}]}
\fi

\usepackage[spanish,es-nodecimaldot]{babel}
\usepackage{microtype}
\usepackage{booktabs}
\usepackage{amssymb}
\usepackage{array}
\usetikzlibrary{arrows.meta,positioning,fit,calc,shapes.geometric}

\graphicspath{{assets/}}

% ---------- Identidad visual UTP ----------
\definecolor{utpRojo}{HTML}{C42B35}
\definecolor{utpNegro}{HTML}{111111}
\definecolor{utpGris}{HTML}{5F5F5F}
\definecolor{utpGrisMedio}{HTML}{9A9A9A}
\definecolor{utpGrisClaro}{HTML}{F1F1F1}
\definecolor{utpNaranja}{HTML}{C1670B}

\usetheme{default}
\setbeamercolor{normal text}{fg=utpNegro,bg=white}
\setbeamercolor{structure}{fg=utpRojo}
\setbeamercolor{frametitle}{fg=utpNegro,bg=white}
\setbeamercolor{title}{fg=utpNegro}
\setbeamercolor{subtitle}{fg=utpNegro}
\setbeamercolor{author}{fg=utpNegro}
\setbeamercolor{date}{fg=utpNegro}
\setbeamercolor{block title}{fg=white,bg=utpNegro}
\setbeamercolor{block body}{fg=utpNegro,bg=utpGrisClaro}
\setbeamercolor{alertblock title}{fg=white,bg=utpRojo}
\setbeamercolor{alertblock body}{fg=utpNegro,bg=utpRojo!7}
\setbeamercolor{exampleblock title}{fg=white,bg=utpGris}
\setbeamercolor{exampleblock body}{fg=utpNegro,bg=utpGris!8}

\setbeamerfont{title}{series=\bfseries,size=\fontsize{19.5}{22.5}\selectfont}
\setbeamerfont{subtitle}{size=\fontsize{11.5}{13.5}\selectfont}
\setbeamerfont{date}{series=\bfseries,size=\fontsize{11.5}{13.5}\selectfont}
\setbeamerfont{author}{size=\fontsize{10.5}{12.5}\selectfont}
\setbeamerfont{frametitle}{series=\bfseries,size=\large}
\setbeamerfont{block title}{series=\bfseries,size=\small}
\setbeamerfont{block body}{size=\normalsize}

\setbeamertemplate{navigation symbols}{}
\setbeamertemplate{headline}{}
\setbeamertemplate{blocks}[default]
\setbeamertemplate{itemize item}{\color{utpRojo}\small\raise.15ex\hbox{$\blacktriangleright$}}
\setbeamertemplate{itemize subitem}{\color{utpGris}\scriptsize$\bullet$}
\setbeamertemplate{enumerate item}{\color{utpRojo}\insertenumlabel.}
\setbeamertemplate{background canvas}{\color{white}\rule{\paperwidth}{\paperheight}}

% Título de diapositiva: negrita negra a la izquierda, logo UTP a la derecha
\setbeamertemplate{frametitle}{%
  \begin{tikzpicture}[remember picture,overlay]
    \node[anchor=north east,inner sep=0pt]
      at ([xshift=-7mm,yshift=-5.5mm]current page.north east)
      {\includegraphics[width=17mm]{utp-logo.png}};
  \end{tikzpicture}%
  \vspace{1.2mm}%
  \begin{minipage}[t]{.80\textwidth}\raggedright\insertframetitle\end{minipage}\par
  \vspace{1.1mm}{\color{utpRojo}\rule{\textwidth}{1.1pt}}
}

% Solo numeración discreta; la plantilla UTP no lleva pie de página
\setbeamertemplate{footline}{%
  \hfill{\color{utpGrisMedio}\scriptsize\insertframenumber}\hspace{6mm}\vspace{2.5mm}%
}

% Portada UTP
\setbeamertemplate{title page}{%
  \begin{tikzpicture}[remember picture,overlay]
    \node[anchor=center,inner sep=0pt] at (current page.center) {%
      \includegraphics[width=\paperwidth,height=\paperheight]{utp-cover-bg.jpeg}%
    };
    \node[anchor=north,align=center,text=utpNegro,inner sep=0pt,
      text width=145mm,font=\usebeamerfont{title}]
      at ([yshift=-25.2mm]current page.north) {\inserttitle};
    \node[anchor=north,align=center,text=utpNegro,inner sep=0pt,
      text width=142mm,font=\usebeamerfont{subtitle}]
      at ([yshift=-38.2mm]current page.north) {\insertsubtitle};
    \node[anchor=north,align=center,text=utpNegro,inner sep=0pt,
      font=\usebeamerfont{date}]
      at ([yshift=-50.4mm]current page.north) {\insertdate};
    \node[anchor=north,align=center,text=utpNegro,inner sep=0pt,
      font=\usebeamerfont{author}]
      at ([yshift=-58.4mm]current page.north) {\insertauthor};
  \end{tikzpicture}%
}

% ---------- Comandos de apoyo didáctico ----------
\newcommand{\concepto}[2]{%
  \begin{block}{#1}\large #2\end{block}}
\newcommand{\pregunta}[1]{%
  \begin{alertblock}{Decide y explica}#1\end{alertblock}}
\newcommand{\analogia}[1]{%
  \begin{exampleblock}{En palabras simples}\small #1\end{exampleblock}}
\newcommand{\enlarealidad}[1]{%
  \begin{block}{Cómo se ve en la realidad}\small #1\end{block}}
\newcommand{\ojo}[1]{%
  \begin{alertblock}{Ojo con esto}\small #1\end{alertblock}}
\newcommand{\apoyo}{{\normalfont\footnotesize\color{utpGris}~· material de apoyo}}
\newcommand{\marca}[1]{\tikz[baseline=(m.base)]{%
  \node[rounded corners=1.5pt,fill=utpGrisClaro,draw=utpGrisMedio,inner xsep=3pt,inner ysep=1.5pt,
        font=\footnotesize] (m) {#1};}}

\tikzset{
  caja/.style={draw=utpNegro, very thick, rounded corners=1.5pt, fill=white,
    align=center, minimum height=0.72cm, inner sep=5pt, font=\small},
  cajaazul/.style={caja, draw=utpGris, fill=white},
  cajaverde/.style={caja, draw=utpNegro, fill=utpGrisClaro},
  cajanaranja/.style={caja, draw=utpRojo!75, fill=utpRojo!7},
  flecha/.style={-{Stealth[length=2.3mm]}, very thick, draw=utpNegro},
  flechavuelta/.style={{Stealth[length=2.3mm]}-{Stealth[length=2.3mm]}, very thick, draw=utpNegro}
}

% ---------- Notación pata de gallo (crow's foot) ----------
% La rotación indica hacia dónde está la entidad: 0 derecha, 180 izquierda,
% 90 arriba, -90 abajo. La pata se abre hacia la entidad del lado "muchos".
\tikzset{
  relacion/.style={thick, draw=utpNegro},
  pics/uno/.style={code={%
    \draw[thick,utpNegro,line cap=round] (0,-1.5mm) -- (0,1.5mm);}},
  pics/muchos/.style={code={%
    \draw[thick,utpNegro,line cap=round,line join=round]
      (2.6mm,1.5mm) -- (0,0) -- (2.6mm,-1.5mm);
    \draw[thick,utpNegro,line cap=round] (0,0) -- (2.6mm,0);}},
}
\newcommand{\simbuno}{\tikz[baseline=-0.6ex]{\draw[relacion](0,0)--(7mm,0);\pic at (4.4mm,0){uno};}}
\newcommand{\simbmuchos}{\tikz[baseline=-0.6ex]{\draw[relacion](0,0)--(7mm,0);\pic at (4.4mm,0){muchos};}}

\usetikzlibrary{decorations.pathreplacing}

\title{¿Qué hay dentro de un gestor?}
\subtitle{Sus piezas, cliente y servidor}
\date{Semana 11 - Sesión 22}
\author{Rel Guzman}

% Palabra reservada de SQL dentro de un bloque de código
\newcommand{\kw}[1]{\textcolor{utpRojo}{\textbf{#1}}}
% Comentario dentro de un bloque de código SQL
\newcommand{\com}[1]{\textcolor{utpGris}{-- #1}}
% Texto que escribe el usuario en la terminal
\newcommand{\tu}[1]{\textcolor{utpRojo}{\textbf{#1}}}
% Insignia de práctica
\newcommand{\ins}[1]{\tikz[baseline=-0.5ex]{\node[rounded corners=2pt,fill=utpRojo,text=white,
  font=\bfseries\scriptsize,inner xsep=5pt,inner ysep=2.5pt] {#1};}}
% Etiqueta oscura: nombre del gestor o del cliente
\newcommand{\etiqueta}[1]{\tikz[baseline=-0.5ex]{\node[rounded corners=2pt,fill=utpNegro,text=white,
  font=\bfseries\scriptsize,inner xsep=6pt,inner ysep=3pt] {#1};}}
% Número de pieza del gestor
\newcommand{\npieza}[1]{\tikz[baseline=-0.6ex]{\node[circle,fill=utpRojo,text=white,
  font=\bfseries\scriptsize,inner sep=1.6pt,minimum size=4.2mm] {#1};}}
% Celda resaltada en rojo claro
\newcommand{\mal}[1]{{\setlength{\fboxsep}{1.2pt}\colorbox{utpRojo!18}{#1}}}
% Barra invertida de las rutas de Windows
\newcommand{\bs}{\textbackslash{}}
% Bloque de código sobre fondo gris
\newsavebox{\sqlbuf}
\newenvironment{sqlcode}[1][.86\textwidth]
  {\begin{center}\begin{lrbox}{\sqlbuf}\begin{minipage}{#1}\ttfamily\footnotesize\vspace{6pt}\raggedright}
  {\vspace{6pt}\end{minipage}\end{lrbox}\colorbox{utpGrisClaro}{\usebox{\sqlbuf}}\end{center}}
% Sangría dentro de un bloque de código
\newcommand{\ind}{\hspace*{1.6em}}
% Anatomía de una instrucción: cada pieza del código con su pregunta al lado
\tikzset{
  pieza/.style={fill=utpGrisClaro,font=\ttfamily\footnotesize,anchor=west,inner sep=5pt,
    minimum height=0.62cm},
  para/.style={font=\small,anchor=west,text=utpNegro}
}
\newcommand{\piezax}{7.6}
\newcommand{\pieza}[3]{%
  \node[pieza] at (0,#1) {#2};
  \node[para] at (\piezax,#1) {\textcolor{utpRojo}{$\blacktriangleleft$}\ #3};}
% Diapositiva de sección: una palabra bien grande (\seccion[0]{...} la deja sin línea roja)
\newcommand{\seccion}[2][1]{%
\begin{frame}[plain]
  \begin{tikzpicture}[remember picture,overlay]
    \node[anchor=north east,inner sep=0pt]
      at ([xshift=-10mm,yshift=-8mm]current page.north east)
      {\includegraphics[width=43mm]{utp-logo.png}};
    \node[anchor=center,inner sep=0pt] at (current page.center)
      {\resizebox{100mm}{!}{\bfseries\color{utpNegro}#2}};
    \ifnum#1=1
    \node[anchor=north,inner sep=0pt] at ([yshift=-11mm]current page.center)
      {\color{utpRojo}\rule{100mm}{1.6pt}};
    \fi
  \end{tikzpicture}
\end{frame}}
% Cabecera de columna dentro de una tabla en monoespaciado
\newcommand{\cab}[1]{\normalfont\bfseries #1}
% Resultado de una consulta: \res{columnas}{cabecera}{filas}
\newcommand{\res}[3]{%
  \centering{\scriptsize\color{utpGris}resultado}\\[3pt]
  \footnotesize\ttfamily
  \begin{tabular}{#1}
  \toprule
  #2 \\
  \midrule
  #3
  \bottomrule
  \end{tabular}}
% Campo de un formulario dibujado: etiqueta y valor
\newcommand{\campo}[3]{%
  \node[anchor=east,font=\scriptsize] at (2.5,#1) {#2};
  \node[anchor=west,draw=utpGrisMedio,fill=white,font=\scriptsize\ttfamily,minimum width=3.6cm,
    minimum height=0.46cm,inner xsep=4pt] at (2.65,#1) {#3};}
% Mensaje de ida (cliente a servidor) y de vuelta en un diagrama de secuencia
\newcommand{\ida}[2]{\draw[flecha] (0,#1)--node[above,font=\tiny,inner sep=1.5pt]{#2}(6.3,#1);}
\newcommand{\vuelta}[2]{\draw[flecha,draw=utpRojo] (6.3,#1)--node[above,font=\tiny,inner sep=1.5pt]{#2}(0,#1);}
% Captura de pantalla con números encima: \captura{ancho}{archivo}{marcadores}
% Dentro de los marcadores, (0,0) es la esquina inferior izquierda y (1,1) la superior derecha.
\newcommand{\captura}[3]{%
  \begin{tikzpicture}
    \node[anchor=south west,inner sep=0pt,draw=utpGrisMedio] (img) at (0,0) {\includegraphics[width=#1]{#2}};
    \begin{scope}[x={(img.south east)},y={(img.north west)}]
      \foreach \px/\py/\n in #3 \node[pin] at (\px,\py) {\n};
    \end{scope}
  \end{tikzpicture}}
\newcommand{\marcaswb}{0.70/0.895/1, 0.66/0.84/2, 0.39/0.706/3, 0.61/0.706/4, 0.60/0.632/5,
  0.463/0.56/6, 0.60/0.481/7, 0.25/0.778/8, 0.63/0.037/9, 0.283/0.037/10}
\newcommand{\marcasssms}{0.30/0.903/1, 0.58/0.597/2, 0.80/0.546/3, 0.80/0.506/4, 0.80/0.468/5,
  0.80/0.355/6, 0.80/0.316/7, 0.72/0.278/8, 0.80/0.241/9, 0.78/0.199/10, 0.80/0.13/11, 0.27/0.037/12}
% Explicación de un campo numerado de una captura
\newcommand{\explica}[3]{\par\noindent\npieza{#1}\ \textbf{#2}\par\vspace{2pt}{\leftskip=6.2mm #3\par}\vspace{10pt}}
% Estilos de los diagramas de esta sesión
\tikzset{
  comp/.style={caja,font=\scriptsize,text width=2.0cm,minimum height=0.8cm,inner sep=3pt},
  grupo/.style={draw=utpGrisMedio,dashed,rounded corners=3pt,inner sep=6pt},
  num/.style={circle,fill=utpRojo,text=white,font=\bfseries\tiny,inner sep=1.2pt,minimum size=3.6mm},
  compu/.style={draw=utpNegro,very thick,rounded corners=4pt,inner sep=8pt},
  vida/.style={utpGrisMedio,thick,dashed},
  pin/.style={circle,fill=utpRojo,text=white,font=\bfseries\scriptsize,inner sep=1pt,
    minimum size=4.4mm,draw=white,thick}
}
\begin{document}
{\setbeamertemplate{footline}{}\maketitle}
\begin{frame}[c]
\begin{block}{Logro de la sesión}
Al finalizar, podrás explicar \textbf{qué piezas tiene un gestor} por dentro, distinguir \textbf{cliente} y \textbf{servidor}, y reconocer en qué se diferencia \textbf{SQLite}, un gestor sin servidor.
\end{block}
\end{frame}
% ---------------------------------------------------------------------
\begin{frame}{Revisión de la sesión anterior}
\small Ejecutaste el modelo físico en SSMS, MySQL Workbench y DB Browser for SQLite.
\vspace{0.4cm}
\begin{center}
\begin{tikzpicture}[font=\small]
  \node[caja,text width=3.3cm,minimum height=1.3cm] (a) {\textbf{Tú}\\{\scriptsize escribes \texttt{CREATE TABLE}\\y presionas \textit{Ejecutar}}};
  \node[caja,cajanaranja,text width=3.3cm,minimum height=1.3cm,right=1.5cm of a] (b) {\Huge\bfseries ?};
  \node[caja,text width=3.3cm,minimum height=1.3cm,right=1.5cm of b] (c) {\textbf{La tabla existe}\\{\scriptsize y mañana sigue ahí}};
  \draw[flecha] (a)--(b); \draw[flecha] (b)--(c);
\end{tikzpicture}
\end{center}
\vspace{0.4cm}
\concepto{La pregunta de hoy}{\small ¿Qué pasa en el signo de interrogación?}
\end{frame}
% ---------------------------------------------------------------------
\begin{frame}{Programa, proceso y servicio}
\begin{description}\setlength\itemsep{10pt}
  \item[Programa] Un archivo con instrucciones. Ejemplos: \texttt{mysqld.exe} (MySQL) y \texttt{sqlservr.exe} (SQL Server).
  \item[Proceso] Un programa \textbf{mientras está funcionando}.
  \item[Servicio] Un proceso que funciona \textbf{sin ventana} y arranca solo al encender Windows.
\end{description}
\vspace{0.3cm}
\small Los servicios se ven en Windows con \textbf{Windows + R} $\rightarrow$ \texttt{services.msc}.
\end{frame}
% ---------------------------------------------------------------------
\begin{frame}{Memoria y disco}
\begin{center}
\renewcommand{\arraystretch}{1.4}
\begin{tabular}{@{}>{\bfseries}l>{\raggedright}p{4.2cm}>{\raggedright\arraybackslash}p{4.2cm}@{}}
\toprule
& \bfseries Memoria (RAM) & \bfseries Disco \\
\midrule
Velocidad & muy rápida & lenta \\
Tamaño & pequeña & grande \\
Al apagar & \mal{se borra} & se conserva \\
\bottomrule
\end{tabular}
\end{center}
\vspace{0.3cm}
\concepto{Por qué importa}{\small Un gestor trabaja en \textbf{memoria} para ser rápido y guarda en \textbf{disco} para no perder nada.}
\end{frame}
% ---------------------------------------------------------------------
\begin{frame}{Base de datos y gestor}
\begin{columns}[c,onlytextwidth]
\column{.56\textwidth}
\begin{description}\setlength\itemsep{10pt}
  \item[Base de datos] Los datos organizados en tablas. Ejemplo: \texttt{peluqueria}.
  \item[Gestor] El \textbf{programa} que guarda las bases de datos y responde pedidos sobre ellas. También se le dice \textbf{SGBD}.
\end{description}
\vspace{0.2cm}
\small Ejemplos: MySQL, SQL Server, PostgreSQL, SQLite.
\column{.4\textwidth}
\centering
\begin{tikzpicture}[font=\small]
  \tikzset{bd/.style={cylinder,shape border rotate=90,aspect=0.22,draw=utpNegro,thick,
    fill=utpGrisClaro,minimum width=1.55cm,minimum height=1.0cm,font=\scriptsize\ttfamily}}
  \node[bd] (b1) at (0,0) {peluqueria};
  \node[bd] (b2) at (1.85,0) {tienda};
  \node[bd] (b3) at (0.92,-1.45) {colegio};
  \node[draw=utpRojo,very thick,rounded corners=4pt,inner sep=9pt,fit=(b1)(b2)(b3),
    label={[font=\small\bfseries]above:Gestor: MySQL}] {};
\end{tikzpicture}
\\[3pt]{\scriptsize Un gestor guarda \textbf{muchas} bases.}
\end{columns}
\end{frame}
% =======================================================================
\seccion{Por dentro}
\begin{frame}{Las piezas de un gestor}
\small Lo que le llega al gestor es una \textbf{consulta}: un pedido escrito en SQL.
\vspace{0.1cm}
\begin{center}
\resizebox{\textwidth}{!}{%
\begin{tikzpicture}[font=\scriptsize]
  \node[caja,cajanaranja,text width=1.7cm] (q) at (0,0) {\textbf{Consulta}};
  \node[comp] (c1) at (3,0) {Gestor de\\conexiones};
  \node[comp] (an) at (6,1.15) {Analizador};
  \node[comp] (op) at (6,0) {Optimizador};
  \node[comp] (ej) at (6,-1.15) {Ejecutor};
  \node[comp] (tr) at (9.1,0.6) {Gestor de\\transacciones};
  \node[comp] (mi) at (9.1,-0.6) {Memoria\\intermedia};
  \node[comp,cajaverde] (lg) at (12.6,1.15) {Registro (log)};
  \node[comp,cajaverde] (ct) at (12.6,0) {Catálogo};
  \node[comp,cajaverde] (ar) at (12.6,-1.15) {Archivos de datos};
  \node[grupo,fit=(an)(op)(ej)] (gp) {};
  \node[font=\scriptsize\bfseries,above=1pt of gp] (lp) {Procesador de consultas};
  \node[grupo,fit=(tr)(mi)] (gm) {};
  \node[font=\scriptsize\bfseries,above=1pt of gm,align=center] (lm) {Motor de\\almacenamiento};
  \node[grupo,fit=(lg)(ct)(ar)] (gd) {};
  \node[font=\scriptsize\bfseries,above=1pt of gd] (ld) {Disco};
  \node[draw=utpRojo,very thick,rounded corners=4pt,inner sep=7pt,fit=(c1)(gp)(lp)(gm)(lm)] (srv) {};
  \node[font=\scriptsize\bfseries,text=utpRojo,below=2pt of srv] {El gestor (en memoria)};
  \draw[flecha] (q)--(c1);
  \draw[flecha] (c1)--(gp.west|-c1);
  \draw[flecha] (an)--(op); \draw[flecha] (op)--(ej);
  \draw[flecha] (gp.east|-0,0)--(gm.west|-0,0);
  \draw[flecha] (gm.east|-0,0)--(gd.west|-0,0);
  \node[num] at (c1.north west) {1};
  \node[num] at (gp.north west) {2};
  \node[num] at (gm.north west) {3};
  \node[num] at (mi.north west) {4};
  \node[num] at (tr.north west) {5};
  \node[num] at (lg.north west) {6};
  \node[num] at (ct.north west) {7};
  \node[num] at (ar.north west) {8};
\end{tikzpicture}}
\end{center}
\vspace{0.1cm}
\small Vamos pieza por pieza.
\end{frame}
% ---------------------------------------------------------------------
\begin{frame}{\npieza{1} Gestor de conexiones}
\concepto{Conexión}{\small Un canal abierto entre quien pide y el gestor. Por él viajan las consultas y las respuestas.}
\vspace{0.2cm}
\begin{itemize}\setlength\itemsep{8pt}
  \item \textbf{Autenticación}: comprueba \textbf{quién eres} (usuario y contraseña).
  \item \textbf{Autorización}: revisa \textbf{qué puedes hacer} (ver, agregar, borrar\ldots).
\end{itemize}
\vspace{0.2cm}
\ojo{El usuario del gestor \textbf{no es} tu usuario de Windows. En MySQL, el administrador se llama \texttt{root}.}
\end{frame}
% ---------------------------------------------------------------------
\begin{frame}{\npieza{2} Procesador de consultas}
\begin{center}
\resizebox{\textwidth}{!}{%
\begin{tikzpicture}[font=\small]
  \tikzset{paso/.style={caja,text width=2.9cm,minimum height=1.4cm}}
  \node[caja,cajanaranja,text width=2.8cm,font=\scriptsize\ttfamily] (q) at (0,0) {SELECT nombre FROM servicio WHERE id = 2;};
  \node[paso] (a) at (3.8,0) {\textbf{Analizador}\\{\scriptsize ¿está bien escrita?\\¿existe la tabla?}};
  \node[paso] (o) at (7.6,0) {\textbf{Optimizador}\\{\scriptsize ¿cuál es el camino más rápido?}};
  \node[paso] (e) at (11.4,0) {\textbf{Ejecutor}\\{\scriptsize sigue el camino y arma el resultado}};
  \draw[flecha] (q)--(a); \draw[flecha] (a)--(o); \draw[flecha] (o)--(e);
\end{tikzpicture}}
\end{center}
\vspace{0.2cm}
\small Si el analizador encuentra un error, se detiene ahí:
\begin{sqlcode}[.8\textwidth]
\scriptsize
\kw{SELECT} nombre \kw{FROM} \mal{servicios};\\
\color{utpRojo}ERROR 1146 (42S02): Table 'peluqueria.servicios' doesn't exist
\end{sqlcode}
\end{frame}
% ---------------------------------------------------------------------
\begin{frame}{\npieza{2} El optimizador y los índices}
\concepto{Índice}{\small Una lista ordenada para encontrar filas rápido, como el índice de un libro. La PK y los UNIQUE ya tienen uno.}
\vspace{0.1cm}
\small \texttt{EXPLAIN} muestra el camino que eligió el optimizador:
\vspace{0.1cm}
\begin{columns}[T,onlytextwidth]
\column{.49\textwidth}
\begin{sqlcode}[\dimexpr\textwidth-2\fboxsep\relax]
\scriptsize
\kw{EXPLAIN SELECT} * \kw{FROM} servicio \kw{WHERE} id = 2;
\end{sqlcode}
\vspace{-0.15cm}
\res{ll}{\cab{type} & \cab{key}}{const & PRIMARY \\}
\\[4pt]\centering\scriptsize usa el índice: va \textbf{directo}
\column{.49\textwidth}
\begin{sqlcode}[\dimexpr\textwidth-2\fboxsep\relax]
\scriptsize
\kw{EXPLAIN SELECT} * \kw{FROM} servicio \kw{WHERE} precio > 40;
\end{sqlcode}
\vspace{-0.15cm}
\res{ll}{\cab{type} & \cab{key}}{\mal{ALL} & NULL \\}
\\[4pt]\centering\scriptsize sin índice: revisa \textbf{todas} las filas
\end{columns}
\end{frame}
% ---------------------------------------------------------------------
\begin{frame}{\npieza{3}\,\npieza{4} Motor de almacenamiento y memoria intermedia}
\begin{columns}[T,onlytextwidth]
\column{.48\textwidth}
\concepto{Motor de almacenamiento}{\small Guarda las filas en el disco y las vuelve a leer. En MySQL es \textbf{InnoDB}: el \texttt{ENGINE = InnoDB} de la sesión anterior.}
\column{.48\textwidth}
\concepto{Memoria intermedia}{\small Una parte de la RAM donde el gestor deja lo que leyó del disco. Si lo vuelve a necesitar, no va al disco.}
\end{columns}
\vspace{0.3cm}
\begin{itemize}\setlength\itemsep{6pt}
  \item Por eso la segunda vez una consulta suele ser \textbf{más rápida}.
  \item Por eso un gestor usa \textbf{mucha RAM}: es a propósito.
\end{itemize}
\end{frame}
% ---------------------------------------------------------------------
\begin{frame}{\npieza{3}\,\npieza{4} ¿Cómo se llaman en cada gestor?}
\begin{center}
\renewcommand{\arraystretch}{1.5}
\begin{tabular}{@{}>{\bfseries}l>{\raggedright}p{3.2cm}>{\raggedright\arraybackslash}p{3.4cm}@{}}
\toprule
Gestor & \bfseries Motor & \bfseries Memoria intermedia \\
\midrule
MySQL & InnoDB & \textit{buffer pool} \\
SQL Server & incluido & \textit{buffer pool} \\
PostgreSQL & incluido & \textit{shared buffers} \\
SQLite & incluido & \textit{page cache} \\
\bottomrule
\end{tabular}
\end{center}
\vspace{0.3cm}
\small Mismo trabajo, otro nombre.
\end{frame}
% ---------------------------------------------------------------------
\begin{frame}{\npieza{5} Gestor de transacciones}
\concepto{Transacción}{\small Un grupo de cambios que ocurren \textbf{todos o ninguno}.}
\vspace{0.05cm}
\begin{sqlcode}[.7\textwidth]
\scriptsize
\kw{START TRANSACTION};\\
\kw{INSERT INTO} pago (monto, metodo, cita\_id) \kw{VALUES} (55.00, 'efectivo', 1);\\
\kw{UPDATE} cita \kw{SET} estado = 'completada' \kw{WHERE} id = 1;\\
\kw{COMMIT}; \ \com{confirma. ROLLBACK deshace}
\end{sqlcode}
\vspace{-0.1cm}
\concepto{Bloqueo}{\small Si dos personas cambian la misma fila a la vez, el gestor hace \textbf{esperar} a la segunda hasta que la primera termine.}
\end{frame}
% ---------------------------------------------------------------------
\begin{frame}{\npieza{6} Registro de recuperación}
\concepto{Registro (\textit{log})}{\small Un archivo donde el gestor \textbf{anota cada cambio} antes de darlo por hecho.}
\vspace{0.2cm}
\begin{center}
\resizebox{.9\textwidth}{!}{%
\begin{tikzpicture}[font=\scriptsize]
  \tikzset{p/.style={caja,text width=2.45cm,minimum height=1.05cm,font=\scriptsize}}
  \node[p] (a) at (0,0) {\textbf{1.} cambia en \textbf{memoria}};
  \node[p,cajanaranja] (b) at (3.3,0) {\textbf{2.} lo anota en el \textbf{registro}};
  \node[p] (c) at (6.6,0) {\textbf{3.} responde «listo»};
  \node[p,cajaverde] (d) at (9.9,0) {\textbf{4.} más tarde, lo guarda en su \textbf{archivo}};
  \foreach \u/\v in {a/b,b/c,c/d} \draw[flecha] (\u)--(\v);
\end{tikzpicture}}
\end{center}
\vspace{0.3cm}
\small Si se va la luz antes del paso 4, al encender el gestor lee el registro y \textbf{recupera} los cambios.
\end{frame}
% ---------------------------------------------------------------------
\begin{frame}{\npieza{7}\,\npieza{8} Catálogo y archivos de datos}
\begin{columns}[T,onlytextwidth]
\column{.48\textwidth}
\concepto{Catálogo}{\small La lista de todo lo que hay: bases, tablas, columnas y usuarios. Son \textbf{datos sobre los datos}.}
\begin{sqlcode}[\dimexpr\textwidth-2\fboxsep\relax]
\scriptsize
\kw{SHOW TABLES};\\
\kw{DESCRIBE} servicio;
\end{sqlcode}
\column{.48\textwidth}
\concepto{Archivos de datos}{\small Donde terminan guardadas las tablas. Cada tabla es un archivo \texttt{.ibd}.}
\begin{sqlcode}[\dimexpr\textwidth-2\fboxsep\relax]
\scriptsize
\kw{SELECT} @@datadir;\\
\com{C:\bs ProgramData\bs MySQL\bs \ldots\bs Data\bs}
\end{sqlcode}
\end{columns}
\vspace{0.1cm}
\ojo{Nunca copies ni borres esos archivos a mano: solo el gestor sabe leerlos.}
\end{frame}
% ---------------------------------------------------------------------
\begin{frame}{Las ocho piezas en una tabla}
\begin{center}
\small
\renewcommand{\arraystretch}{1.3}
\begin{tabular}{@{}c>{\raggedright\bfseries}p{4.3cm}>{\raggedright\arraybackslash}p{6.8cm}@{}}
\toprule
& Pieza & \bfseries Para qué sirve \\
\midrule
\npieza{1} & Gestor de conexiones & deja entrar solo a quien corresponde \\
\npieza{2} & Procesador de consultas & entiende el SQL y busca el camino más rápido \\
\npieza{3} & Motor de almacenamiento & guarda y lee las filas \\
\npieza{4} & Memoria intermedia & evita ir al disco cada vez \\
\npieza{5} & Gestor de transacciones & todo o nada; uno a la vez \\
\npieza{6} & Registro & no perder nada si se va la luz \\
\npieza{7} & Catálogo & saber qué existe \\
\npieza{8} & Archivos de datos & guardar todo en el disco \\
\bottomrule
\end{tabular}
\end{center}
\end{frame}
% =======================================================================
\seccion[0]{Cliente y servidor}
\begin{frame}{Cliente y servidor}
\begin{columns}[T,onlytextwidth]
\column{.48\textwidth}
\concepto{Servidor}{\small El programa que \textbf{espera pedidos} y los atiende. El gestor es un servidor.}
\column{.48\textwidth}
\concepto{Cliente}{\small El programa que \textbf{envía pedidos} al servidor y muestra las respuestas.}
\end{columns}
\vspace{0.3cm}
\begin{center}
\begin{tikzpicture}[font=\scriptsize]
  \node[caja,text width=2.5cm] (c1) at (0,1.1) {Terminal};
  \node[caja,text width=2.5cm] (c2) at (0,0) {MySQL Workbench};
  \node[caja,text width=2.5cm] (c3) at (0,-1.1) {App de la peluquería};
  \node[caja,cajanaranja,text width=2.2cm,minimum height=1.3cm] (s) at (4.6,0) {\textbf{Servidor}\\MySQL};
  \foreach \c in {c1,c2,c3} \draw[flechavuelta] (\c.east)--(s.west);
  \node[font=\scriptsize\bfseries,above=1mm of c1] {Clientes};
\end{tikzpicture}
\end{center}
\vspace{0.1cm}
\small Muchos clientes, un servidor. Solo el servidor toca los datos.
\end{frame}
% ---------------------------------------------------------------------
\begin{frame}{¿Quién es quién?}
\begin{center}
\renewcommand{\arraystretch}{1.4}
\begin{tabular}{@{}>{\bfseries}l>{\raggedright}p{3.9cm}>{\raggedright\arraybackslash}p{4.6cm}@{}}
\toprule
Gestor & \bfseries Servidor & \bfseries Clientes \\
\midrule
MySQL & \texttt{mysqld.exe}\\{\footnotesize servicio \texttt{MySQL80}} & terminal, MySQL Workbench \\
SQL Server & \texttt{sqlservr.exe}\\{\footnotesize servicio \texttt{SQL Server (\ldots)}} & SSMS \\
\bottomrule
\end{tabular}
\end{center}
\vspace{0.3cm}
\concepto{Workbench y SSMS no son la base de datos}{\small Son \textbf{clientes}. Si cierras Workbench, los datos siguen en el servidor. Si apagas el servicio, ningún cliente puede entrar.}
\end{frame}
% ---------------------------------------------------------------------
\begin{frame}{SQLite: un gestor sin servidor}
\begin{center}
\renewcommand{\arraystretch}{1.4}
\begin{tabular}{@{}>{\raggedright\bfseries}p{3.4cm}>{\raggedright}p{4.0cm}>{\raggedright\arraybackslash}p{3.6cm}@{}}
\toprule
& \etiqueta{MySQL, SQL Server} & \etiqueta{SQLite} \\
\midrule
¿Hay servidor? & sí & no \\
¿Dónde está la base? & en el servidor & en un archivo \texttt{.db} \\
¿Cómo se entra? & conectándose & abriendo el archivo \\
\bottomrule
\end{tabular}
\end{center}
\vspace{0.3cm}
\small Por eso DB Browser for SQLite no pide usuario ni contraseña. La próxima sesión veremos cómo conectarse a un gestor \textbf{con servidor}.
\end{frame}
\begin{frame}[plain]
  \begin{tikzpicture}[remember picture,overlay]
    \node[anchor=north west,font=\bfseries\fontsize{24}{28}\selectfont]
      at ([xshift=12mm,yshift=-9mm]current page.north west) {Preguntas};
    \node[anchor=north east,inner sep=0pt]
      at ([xshift=-10mm,yshift=-8mm]current page.north east)
      {\includegraphics[width=43mm]{utp-logo.png}};
    \node[anchor=center,inner sep=0pt]
      at ([yshift=-5mm]current page.center)
      {\includegraphics[width=78mm]{preguntas-ilustracion.png}};
  \end{tikzpicture}
\end{frame}
\begin{frame}{Cierre}
\begin{center}
\renewcommand{\arraystretch}{1.4}
\begin{tabular}{>{\raggedright\bfseries}p{3.4cm}>{\raggedright\arraybackslash}p{8.4cm}}
\toprule
Gestor & El programa que guarda las bases de datos. \\
Servidor & El programa que atiende; ahí están los datos. \\
Cliente & El programa que pide: terminal, Workbench, SSMS. \\
Piezas & Conexiones, consultas, almacenamiento, memoria, transacciones, registro, catálogo y archivos. \\
SQLite & Un gestor sin servidor: la base es un archivo. \\
\bottomrule
\end{tabular}
\end{center}
\end{frame}

\begin{frame}{Referencias}
\footnotesize
\begin{itemize}\setlength\itemsep{8pt}
\item Silberschatz, A., Korth, H. F., \& Sudarshan, S. (2014). \textit{Fundamentos de bases de datos} (6.ª ed.). McGraw-Hill.
\item Oracle. MySQL Reference Manual: \href{https://dev.mysql.com/doc/refman/8.0/en/innodb-buffer-pool.html}{\textit{Buffer Pool}} y \href{https://dev.mysql.com/doc/refman/8.0/en/innodb-redo-log.html}{\textit{Redo Log}}.
\item SQLite. \href{https://www.sqlite.org/serverless.html}{\textit{SQLite Is Serverless}}.
\end{itemize}
\end{frame}
% Cierre UTP
{
  \setbeamertemplate{background canvas}{%
    \includegraphics[width=\paperwidth,height=\paperheight]{gracias-final.pdf}%
  }
  \setbeamertemplate{footline}{}
  \begin{frame}[plain]
  \end{frame}
}

\end{document}
